% Generated draft: requires mathematical and source review.
\documentclass[11pt]{article}
\usepackage[margin=25mm]{geometry}
\usepackage{fontspec}
\setmainfont{Latin Modern Roman}
\usepackage{amsmath,amssymb,mathtools,unicode-math}
\setmathfont{Latin Modern Math}
\usepackage{enumitem,needspace,xurl,hyperref}
\setlength{\parindent}{0pt}
\setlength{\parskip}{5pt}
\setlength{\emergencystretch}{3em}
\newcommand{\N}{\mathbb{N}}
\newcommand{\Z}{\mathbb{Z}}
\newcommand{\Q}{\mathbb{Q}}
\newcommand{\R}{\mathbb{R}}
\newcommand{\C}{\mathbb{C}}
\newcommand{\F}{\mathbb{F}}
\newcommand{\A}{\mathbb{A}}
\newcommand{\PP}{\mathbb{P}}
\newcommand{\E}{\mathbb{E}}
\newcommand{\eps}{\varepsilon}
\newcommand{\dd}{\,\mathrm{d}}
\begin{document}
\section*{UnsolvedMath ID: 30003515}
\textbf{Source problem code:} catalogue entry
\subsection{Definitions and mathematical statement}
Let \(F\) denote the field appearing in the supplied statement, and write \(\operatorname{GL}_6(F)\) for the group of invertible \(6\times6\) matrices over \(F\). The record discusses archimedean and \(p\)-adic cases but does not specify the precise allowed fields, their characteristics, or the category of representations.

Let \(\pi\) be a representation described by the source as generic, with central character \(\omega_\pi=\chi^2\). Under the usual interpretation of this hypothesis, \(\chi\) is a character of \(F^\times\), and
\[
\omega_\pi(a)=\chi(a)^2\qquad(a\in F^\times).
\]
The source does not explicitly state irreducibility, admissibility, continuity or smoothness requirements, or a definition of genericity.

Retain the source notation
\[
E(\pi,\chi):=\epsilon\!\left(\frac12,(\wedge^3\pi)\otimes\chi^{-1}\right).
\]
Here \(\wedge^3\pi\) is the source's exterior-cube notation; its precise construction and the conventions for the epsilon factor are not supplied. Likewise, \(m(\pi)\) is not defined in the record. The cited Ginzburg--Rallis reference suggests a model-multiplicity interpretation, but the relevant subgroup, character, and Hom space are not specified and are not reconstructed here.

\textbf{Source-stated epsilon-dichotomy conjecture.} For every generic \(\pi\) and character \(\chi\) satisfying the central-character hypothesis, within the intended but incompletely specified setting, do both equivalences hold?
\[
\begin{aligned}
m(\pi)=1&\quad\Longleftrightarrow\quad E(\pi,\chi)=1,\\
m(\pi)=0&\quad\Longleftrightarrow\quad E(\pi,\chi)=-1.
\end{aligned}
\]
Both equivalences are retained as requested assertions. Calling the second equivalent to the first requires additional range information, such as \(m(\pi)\in\{0,1\}\) and \(E(\pi,\chi)\in\{1,-1\}\), which the supplied record does not establish. Until the field, representation category, multiplicity, and epsilon-factor conventions are specified, this is a qualified transcription of the proposed target rather than a fully determined mathematical conjecture.
\subsection{Short English statement}
The supplied conjecture predicts that a generic representation of \(\operatorname{GL}_6(F)\) with central character \(\chi^2\) has multiplicity one exactly when its twisted exterior-cube epsilon factor at \(1/2\) is \(+1\), and multiplicity zero exactly when that factor is \(-1\). The intended multiplicity and several hypotheses and conventions remain unspecified. The record labels the problem partially solved; that assessment is supplied dated curation requiring independent verification.
\subsection{Sources}
\begin{itemize}
\item Supplied local UnsolvedMath registry record: permanent ID \(30003515\), title ``Epsilon Dichotomy for Generic General Linear Representations,'' source problem code \texttt{catalogue entry}, in the Oberwolfach Reports Open Problems collection. The supplied proposed year is 2017; no individual proposer is identified.
\item Supplied report citation: \emph{Harmonic Analysis and the Trace Formula} (2018), \emph{Oberwolfach Reports}, pp.\ 1551--1630, DOI \texttt{10.4171/owr/2017/25}. \url{https://doi.org/10.4171/owr/2017/25}. The record describes these reports as documentation of mathematical workshops and their open-problem discussions. No specific page or named contribution within the report is supplied.
\item Supplied literature source, labeled ``Wan, local Ginzburg--Rallis model'': \url{https://sites.rutgers.edu/chen-wan/wp-content/uploads/sites/518/2020/08/GRgeneric.pdf}.
\item Supplied literature source, labeled ``local trace formula for the generalized Shalika model'': \url{https://hal.science/hal-01941350v1/document}.
\item Supplied literature source, labeled ``original OWR report'': \url{https://doi.org/10.4171/owr/2017/25}.
\item Literature assessment dated exactly \texttt{2026-08-21}: the supplied curation labels the problem partially solved and reports proofs in archimedean cases and broad \(p\)-adic tempered or non-discrete classes, with exceptional \(p\)-adic discrete-series cases and the full generic case remaining unresolved. These are supplied dated curation claims, not independently verified mathematical or literature conclusions. The record itself requests independent verification of time-sensitive status and cited claims.
\item The supplied statement status is \texttt{corrected\_verified}, with an editorial explanation that mathematical content was preserved while notation, typography, and sentence structure were repaired. This supplied label does not constitute independent mathematical review of this draft or verification of the cited sources.
\end{itemize}
\end{document}
